\documentclass[instructions.tex]{subfiles}
\begin{document}
 
\chapter{L'espérance du paradis nous console et nous encourage à servir Dieu.}
 
\primo Il est consolant de travailler pour un maître qui sait récompenser abondamment. Dieu, le meilleur de tous les maîtres, récompense avec tant de magnificence les moindres services, que rien ne se perd de tout ce que nous faisons pour sa gloire : il nous tient même compte de nos bons désirs. Le moindre degré du bonheur qu'il promet à l'homme juste, est infiniment plus estimable que tout l'or de la terre.
 
Tout ce que nous faisons pour Dieu est comme un dépôt, un trésor que nous lui confions. Un homme qui ouvre un trésor qu'il a amassé pendant vingt ans, retrouve dans ce trésor tout ce qu'il y a mis, jusqu'aux plus petites pièces de monnoie. De même, à notre mort, nous trouverons dans le trésor de la bonté de Dieu tout ce que nous aurons fait pour lui. La moindre action faite en état de grâce pour sa gloire, ne fût-ce qu'un verre d'eau donné pour l'amour de Jésus-Christ, ou un acte de patience, aura une récompense au-dessus de tout ce que nous pouvons comprendre. Il nous demande peu, et nous rendra beaucoup, dit saint Pierre Chrysologue\footnote{Parva petit, maxima reddituras. S. Chrys.}. 
 
Quelle perte ne faisons-nous pas, lorsque nous négligeons les occasions de vertus qui se présentent chaque jour! Si les saints dans le ciel étaient capables de tristesse, ils gémiraient d'avoir laissé échapper tant d'occasions de plaire à Dieu, et d'augmenter, par les plus petites choses, leur mérite et leur gloire. Consolez-vous, chrétiens, vous pouvez, dans vos disgrâces, dans vos travaux et dans vos occupations, remplir vos années de mérites pour le ciel, par la résistance aux tentations, par de fréquens retours de votre cœur à Dieu, par des actes de patience, de charité, d'obéissance, etc.
 
Si vous étiez pénétré de ces vérités, les œuvres de piété, l'aumône, la mortification, seraient vos délices : votre âme s'élèverait souvent à Dieu pour l'aimer et l'adorer : plein de charité et de zèle, vous seriez en paix avec tous; vous n'auriez point de plus douce consolation que de vivre dans l'espérance de posséder bientôt votre Dieu dans sa gloire. Que les choses de la terre me paraissent méprisables, lorsque je contemple la gloire du ciel, disait saint Ignace\footnote{Quàm sordet terra, dùm cœlum aspicis! S. Ign.}.
 
Autant l'espérance du ciel nous console, autant elle nous remplit de courage. Si la peine qu'il faut prendre pour le gagner vous effraie, dit saint Augustin, que la récompense vous anime\footnote{Si labor terret, merces invitet. S. Aug.}. Cette récompense n'est pas éloignée : bientôt vos travaux finiront, mais votre bonheur et votre joie ne finiront point : Je travaille pour le ciel, disait un martyr; cette pensée vous remplira de force. Les saints vous crient du haut du ciel : Oh! que vous êtes heureux sur la terre, de pouvoir gagner le bonheur dont nous jouissons! Les âmes dans les feux du purgatoire vous crient : Oh! que vous êtes heureux de pouvoir éviter ces flammes pour aller bientôt jouir de Dieu! Les damnés dans leurs tourmens vous crient : Oh! que vous êtes heureux de pouvoir gagner le ciel que nous avons perdu! Ah! si nous avions le temps que vous avez! N'en perdez pas un moment; faites-vous sages à nos dépens. Votre âme elle-même vous crie : Oh! que vous êtes malheureux de m'abandonner, de ne rien faire pour me procurer le ciel, tandis qu'on se donne tant de peines pour des vanités!
 
En effet, un soldat ne craint point les fatigues, quand il s'agit du service de son prince; il va au combat pour acquérir une fumée de gloire : pourquoi vous épargneriez-vous, quand il s'agit de servir Dieu, et d'acquérir une gloire éternelle\footnote{Illi quidem ut corruptibilem coronam accipiant, nos autem incorruptam. Cor. 9. 25.}. Le marchand travaille, s'expose à mille dangers pour faire une fortune dont il ne jouira que peu de temps, ou dont il ne jouira même jamais : pourquoi vous plaindriez-vous en travaillant pour acquérir dans le ciel un bonheur sans fin? Oh! que nous sommes ingrats et misérables, si nous nous épargnons en servant un Dieu qui nous aime, qui nous récompense si libéralement! Quelques momens de tribulation seront payés par une gloire immense\footnote{Momentaneum et leve tribulationis nostræ suprà modum in sublimitate æternum gloriæ pondus operatur in nobis. 2. Cor. 4.}.
 
\section{Résolutions}
 
\primo Je me rappellerai souvent cette pensée : Je puis encore prétendre au ciel.
 
\secundo Je considérerai aussi que si j'ai le bonheur d'être en grâce avec Dieu, de faire mes actions, d'endurer mes souffrances, de pratiquer les devoirs de mon état, d'exercer des actes de vertus, en vue de plaire à Dieu, je mériterai par tous ces efforts une récompense éternelle.
 
\tertio Je ne craindrai rien tant que de perdre l'amitié de Dieu par le péché mortel, et le fruit de mes œuvres, en les faisant par des motifs humains ou mauvais.
 
\prayer{Que votre bonté pour nous est grande, ô mon Dieu! puisque vous récompensez si largement ceux qui vous aiment, et qui vous servent avec fidélité !}
 
\end{document}
